% !TeX program = pdflatex
% papers/sql.tex — O campo de contagem G_q entre a tabela e a arena:
%   sql_hist_jev do tiffany-db, nos tres niveis (matematica / contrato SQL / C).
%
% Fontes da escrita:
%   tiffany-db/banco/sql.c        (sql_hist_jev:16261, spec_decode:16233,
%                                   celula_valor:7260, S_VIVO/S_PRES, INSERT,
%                                   DELETE, UPDATE)
%   tiffany-db/banco/sql_api.h:148
%   tiffany-db/banco/jev.c        (jev_proj_raw:85, jev_marginal_raw:97,
%                                   jev_escore_raw:156)
%   tiffany-db/banco/jev_api.h, jev_core.h
%   tiffany-db/tests/test_sql_hist_jev.c
%   tiffany-db/FASE_4_9_R3_RESULTADO.md
%
% Auto-contido: cd papers && pdflatex sql.tex
\documentclass[11pt,a4paper]{article}

\usepackage[utf8]{inputenc}
\usepackage[T1]{fontenc}
\usepackage[portuguese]{babel}

\usepackage{amsmath,amssymb,amsthm}

\usepackage[margin=2.4cm]{geometry}

\usepackage{booktabs}
\usepackage{array}

\usepackage{xcolor}

\usepackage[hidelinks,breaklinks]{hyperref}

\sloppy
\emergencystretch=2em

\newtheorem{teorema}{Teorema}
\newtheorem{proposicao}[teorema]{Proposi\c{c}\~ao}
\newtheorem{corolario}[teorema]{Corol\'ario}
\newtheorem{lema}[teorema]{Lema}
\newtheorem{definicao}{Defini\c{c}\~ao}
\newtheorem{observacao}{Observa\c{c}\~ao}

\newcommand{\code}[1]{\texttt{\small #1}}
\newcommand{\cp}[1]{\texttt{\detokenize{#1}}}
\newcommand{\dg}{^{\dagger}}
\newcommand{\Word}{\mathcal{W}}
\newcommand{\FF}{\mathbb{F}}
\newcommand{\Qq}{\mathbb{Q}}
\newcommand{\N}{\mathbb{N}}
\newcommand{\Zz}{\mathbb{Z}}
\newcommand{\Un}[1]{\underline{#1}}
\newcommand{\nr}{\mathrm{nr}}
\newcommand{\nc}{\mathrm{nc}}
\newcommand{\ind}[1]{\mathbf{1}_{\{#1\}}}
\newcommand{\sti}{\textbf{Implementado e demonstrado.}}
\newcommand{\stipa}{\textbf{Implementado, mas interpreta\c{c}\~ao matem\'atica
  parcial.}}
\newcommand{\stco}{\textbf{Compat\'ivel, mas n\~ao demonstrado.}}
\newcommand{\stni}{\textbf{N\~ao estabelecido pela implementa\c{c}\~ao
  atual.}}

\definecolor{cinza}{gray}{0.35}
\newcommand{\medido}[1]{\par\medskip\noindent{\small\color{cinza}\textbf{Medido.} #1}\par\medskip}

% ─────────────────────────────────────────────────────────────────────────────
%  papers/gkcapa.tex — a capa do Reino, mínima, para os papers em `article`.
%
%  O estilo.tex completo é do `report` (títulos de capítulo, skak, …). Aqui fica
%  só o que a capa precisa, para o paper continuar a compilar sozinho com
%  pdflatex. O tradutor próprio (tests/tex) carrega o estilo.tex da raiz / o
%  symlink papers/estilo.tex e avalia o mesmo `\gkcapa`.
% ─────────────────────────────────────────────────────────────────────────────
\definecolor{ouro}{HTML}{B8912F}
\definecolor{ouroclaro}{HTML}{F3E9CE}
\definecolor{tinta}{HTML}{1E1B16}
\definecolor{regua}{HTML}{6E6858}
\providecommand{\gknota}{\fontsize{7.62}{11.05}\selectfont}
\providecommand{\gktexto}{\fontsize{10.50}{15.22}\selectfont}
\providecommand{\gksub}{\fontsize{12.33}{17.87}\selectfont}
\providecommand{\gksec}{\fontsize{14.47}{20.98}\selectfont}
\providecommand{\gkcap}{\fontsize{16.99}{24.63}\selectfont}
\providecommand{\gktit}{\fontsize{23.42}{33.95}\selectfont}
\providecommand{\Prj}{\mathbb{P}}
\providecommand{\Trf}{\mathcal{T}}
\providecommand{\gkcapa}[3]{%
\title{\vspace{-0.6cm}
{\color{ouro}\rule{\textwidth}{1.4pt}}\\[6mm]
\textbf{\gktit\color{tinta} Reino Dourado}\\[3mm]{\gksec\itshape\color{regua} Golden Kingdom}\\[8mm]
{\gkcap\scshape\color{ouro} #1}\\[5mm]
{\gksub\color{regua} #2}\\[2mm]
{\gktexto\color{regua}\itshape #3}\\[7mm]
{\color{ouro}\rule{\textwidth}{0.6pt}}\\[9mm]{\gknota\color{regua}\begin{minipage}{\textwidth}\setlength{\parindent}{0pt}%
\textbf{Aviso de Isenção Algorítmica:} O universo, as leis e as entidades contidas nestas páginas — incluindo felinos matriciais, aves de rapina algébricas e amuletos de controle de fluxo — são construções de ficção formal. Esta obra constrói objetos; não reivindica, por si, descrever a realidade física, fazer engenharia reversa do cosmos ou estabelecer doutrina religiosa. Qualquer coincidência com o universo é assunto para medição; qualquer semelhança com bugs reais do seu sistema operacional é mera coincidência (ou falta de reza). Prossiga por sua própria conta e risco.\end{minipage}}}
\author{}\date{}}


\begin{document}
\gkcapa{O campo de contagem $G_q$ entre a tabela e a arena}
       {$G_q(j)=\displaystyle\sum_{x\in q^{-1}(j)}G(x)$,
        com $G(x)=\lvert \pi^{-1}(x)\rvert$, no caminho C do tiffany-db}
       {sql\_hist\_jev, spec base-36, jev\_marginal, S\_VIVO / S\_PRES /
        S\_LINHAS / S\_ALTO, celula\_valor, guardas E1..E7 e F17}
\maketitle

%=======================================================================
\begin{abstract}
Este paper documenta o que a fun\c{c}\~ao \cp{sql_hist_jev} do
\cp{tiffany-db} realmente faz: l\^e uma coluna de uma tabela, conta as
linhas vivas por valor, constr\'oi o campo de contagem $G$, aplica a
parti\c{c}\~ao por fronteiras $q$ dada pela \cp{spec} base-36, e devolve a
marginal $G_q(j)=\sum_{x\in q^{-1}(j)}G(x)$ com a soma de conserva\c{c}\~ao
$\sum_j G_q(j)=\lvert I\rvert$. O paper cobre tamb\'em a camada relacional que
o motor implementa --- selec\c{c}\~ao, projec\c{c}\~ao, agrupamento, jun\c{c}\~ao,
ordem e escritas ---, modelada como opera\c{c}\~oes matem\'aticas sobre a tabela
ao longo da cadeia
$\text{SQL}\to\text{objecto matem\'atico}\to\text{representa\c{c}\~ao
Tiffany}\to\text{JEV, quando aplic\'avel}$ (Sec.~\ref{sec:relacional}). Cada afirma\c{c}\~ao \'e marcada com um de
tr\^es carimbos: \textbf{Contrato implementado.} (o que a interface diz),
\textbf{Interpreta\c{c}\~ao matem\'atica da implementa\c{c}\~ao.} (o objecto
matem\'atico que o c\'odigo sustenta), e \textbf{N\~ao estabelecido pela
implementa\c{c}\~ao atual.} (o que o c\'odigo n\~ao faz). N\~ao se importa
teoria: o que n\~ao est\'a no c\'odigo n\~ao \'e afirmado.
\end{abstract}

%=======================================================================
\section{O que este paper mede}

\paragraph{Tr\^es n\'iveis.} Todo o enunciado fica num de tr\^es n\'iveis:
\begin{itemize}
\item \textbf{N\'ivel 1 (Matem\'atica):} objectos — campo de contagem,
  parti\c{c}\~ao, marginal — e as suas identidades ($\sum_x G(x)=\lvert I\rvert$,
  $\sum_j G_q(j)=\lvert I\rvert$).
\item \textbf{N\'ivel 2 (Contrato SQL):} a assinatura p\'ublica
  \cp{sql_hist_jev}, o que ela recebe e o que devolve em \cp{SqlOut};
  o c\'odigo de erro em \cp{out->err}.
\item \textbf{N\'ivel 3 (Implementa\c{c}\~ao C):} mem\'oria e datas —
  \cp{mem_le}, \cp{celula_valor}, os bitmaps \cp{S_VIVO}/\cp{S_PRES},
  o topo \cp{cat_nrows}, o layout da arena, os offsets $D_1,J,D_2,\mathrm{OUT}$.
\end{itemize}

\paragraph{A cadeia.} Quando a opera\c{c}\~ao se aplica \`a coluna-alvo da
leitura JEV, o caminho completo \'e
\[
\text{SQL}
\;\longrightarrow\;
\text{opera\c{c}\~ao matem\'atica}
\;\longrightarrow\;
\text{representa\c{c}\~ao Tiffany}
\;\longrightarrow\;
\text{JEV},
\]
em que o \'ultimo passo \'e o que as Secs.~\ref{sec:nivel1}--\ref{sec:porta}
documentam. A Sec.~\ref{sec:relacional} preenche os passos anteriores: \'e uma
camada \emph{anterior e complementar} \`a JEV, e n\~ao a substitui.

\paragraph{Carimbos.} \textbf{Contrato implementado.} e
\textbf{Interpreta\c{c}\~ao matem\'atica da implementa\c{c}\~ao.} acompanham
afirma\c{c}\~oes verific\'aveis no c\'odigo; \textbf{N\~ao estabelecido pela
implementa\c{c}\~ao atual.} marca aquilo que este ficheiro n\~ao faz e que,
portanto, este paper n\~ao atribui.

\paragraph{O que n\~ao se afirma.} N\~ao se afirma aqui que o motor SQL
execute a \`{a}lgebra do \cp{jev.tex}: o caminho que ele executa n\~ao calcula
nem $p_q=G_q/\lvert I\rvert$, nem o funcional $L$, nem probabilitate. O que o
c\'odigo calcula \'e a marginal $G_q$ — nada mais — e \'e isso que se documenta.

%=======================================================================
\section{N\'ivel 2 — o contrato SQL}

\begin{definicao}[Objeto]
\label{def:sqlhistjev}
\textbf{Contrato implementado.}
\begin{center}
\ttfamily\small
\detokenize{int sql_hist_jev(const char *tabela, const char *coluna,}\\
\detokenize{                     SqlOut *out, long *pi1, long *pi2,}\\
\detokenize{                     const char *spec)}
\end{center}
declarada em \cp{banco/sql_api.h:148}, definida em \cp{banco/sql.c:16261}
(n\~ao \'e \cp{static}). Devolve $1$ em sucesso e $0$ em recusa, com a
mensagem em \cp{out->err} quando \cp{out} existe. Os par\^ametros obrigat\'orios
s\~ao \cp{tabela}, \cp{coluna}, \cp{out}, \cp{pi1} (o \'unico canal de
$\lvert I\rvert$) e \cp{spec}. O par\^ametro \cp{pi2} \'e ignorado
(\cp{sql.c:16273-16274}): as duas dimens\~oes da conjunta s\~ao
\cp{c1 = c2 = 1} por decis\~ao F10, n\~ao lidas de \cp{pi2}.
\end{definicao}

\begin{observacao}[Alcance]
\textbf{Contrato implementado.} \cp{sql_hist_jev} n\~ao \'e alcan\c{c}avel por
texto SQL: \cp{sql.c} s\'o chama \cp{jev_arena} e \cp{jev_marginal} dentro
dela; n\~ao existe sintaxe SQL que a invoque. O acesso \'e pela porta C
(declarada em \cp{sql_api.h}) — testes e \cp{harness} s\~ao os consumidores.
\end{observacao}

\begin{definicao}[Sa\'ida]
\label{def:saida}
\textbf{Contrato implementado.} Em sucesso,
\cp{out->nrows = k} (o n\'umero de intervalos), \cp{out->ncols = 3}, e as tr\^es
colunas s\~ao de tipo \cp{SQL_TIPO_INT4}:
\cp{classe} ($j\in\cp{0..k-1}$), \cp{gq} ($G_q(j)$), \cp{i_total}
($\lvert I\rvert$, o mesmo valor devolvido por \cp{*pi1}). As c\'elulas
\cp{nulo[][]} s\~ao zeradas e não h\'a c\'elula ausente na sa\'ida
(\cp{sql.c:16472-16483}).
\end{definicao}

\textbf{Contrato implementado.} \'Os erros E1..E7 (Sec.~\ref{sec:guardas})
devolvem $0$ e gravam a raz\~ao em \cp{out->err}; para além dos E-c\'odigos
existem as recusas \cp{argumentos nulos}, \cp{jev_marginal falhou} e
\cp{conservacao violada}. Em toda a recusa a tabela anterior em
\cp{g_tabela} \'e restaurada.

%=======================================================================
\section{N\'ivel 1 — a tabela como objecto matem\'atico}
\label{sec:tabela}

\begin{definicao}[Dom\'inios de \'indice]
Escreve-se $\Un{n}=\{0,1,\ldots,n-1\}$. O dom\'inio dos valores de uma c\'elula
\cp{INTEIRO} \'e $V=\{0,\ldots,65535\}$ (16 bits sem sinal).
\end{definicao}

\begin{definicao}[Tabela]
\label{def:tabela}
\textbf{Interpreta\c{c}\~ao matem\'atica da implementa\c{c}\~ao.} Uma tabela
\'e um qu\'adruplo
\[
T=(I,J,\mathrm{pres},v),
\]
com: $I\subseteq\Un{nr}$ o conjunto de \emph{linhas vivas}; $J=\Un{nc}$ o
conjunto de colunas; $\mathrm{pres}\colon I\times J\to\{0,1\}$ a \emph{presen\c
{c}a} por c\'elula; e $v\colon\{(i,j)\in I\times J:\mathrm{pres}(i,j)=1\}\to V$
o \emph{valor} das c\'elulas presentes. A linha $i$ \'e a fun\c{c}\~ao parcial
$r_i\colon J\dashrightarrow V$ dada por $r_i(j)=v(i,j)$ se
$\mathrm{pres}(i,j)=1$, e indefinida caso contr\'ario.
\end{definicao}

\begin{observacao}[Porque uma fam\'ilia de parciais, e n\~ao uma rela\c{c}\~ao]
\textbf{Interpreta\c{c}\~ao matem\'atica da implementa\c{c}\~ao.} O c\'odigo
sustenta esta abstrac\c{c}\~ao e n\~ao as alternativas:
\begin{itemize}
\item \emph{n\~ao} \'e uma rela\c{c}\~ao $R\subseteq X\times Y$: n\~ao h\'a
  chave de valor; valores repetem-se entre linhas, e as linhas distinguem-se
  pelo \'indice $i$, n\~ao pelo conte\'udo;
\item a c\'elula \'e endere\c{c}ada pelo produto $i\cdot nc+j$
  (\cp{cel_ix(i*nc+j)}, \cp{sql.c:895}), o que faz de $(I\times J)$ o gr\'afico
  de uma fam\'ilia indexada e n\~ao de um atributo de um produto;
\item a aus\^encia \'e um bit pr\'oprio por c\'elula (\cp{S_PRES}), n\~ao um
  valor convencionado — logo $r_i$ \'e realmente parcial;
\item o bit \cp{S_VIVO} separa \emph{existe a linha} de \emph{n\~ao existe},
  o que explica $I\subseteq\Un{nr}$ (o \'indice \'e topo, n\~ao cardinal).
\end{itemize}
\end{observacao}

\begin{observacao}[O que \'e e o que n\~ao \'e um nulo]
\textbf{Interpreta\c{c}\~ao matem\'atica da implementa\c{c}\~ao.}
\begin{itemize}
\item \textbf{Aus\^encia:} $\mathrm{pres}(i,j)=0$ — a c\'elula n\~ao tem
  valor. Nasce assim por omiss\~ao num \cp{INSERT INTO t (a) VALUES (\dots)}
  em que a coluna n\~ao \'e nomeada (\cp{sql.c:3795}, \emph{``nascem
  AUSENTES''}), a menos que um \cp{DEFAULT} ocupe o seu lugar
  (\cp{sql.c:3800-3801}). Um \cp{NULL} expl\'icito tamb\'em grava
  aus\^encia (\cp{nulo[j]=1}, \cp{sql.c:3689-3690}).
\item \textbf{Escrever faz existir:} um valor posto \'e uma presen\c{c}a — o
  \cp{UPDATE} acende \cp{S_PRES} nas c\'elulas que tocou
  (\cp{sql.c:6126-6133}); mas \cp{SET col = NULL} \emph{desliga} a presen\c{c}a
  (\cp{set_anula ? 0 : 1}).
\item \textbf{Zero \'e um valor:} $\mathrm{pres}(i,j)=1$ com $v(i,j)=0$ conta
  no campo em $G(0)$; n\~ao \'e aus\^encia.
\item \textbf{Linha morta \'e linha fora de $I$:} um \cp{DELETE} s\'o desliga o
  bit \cp{S_VIVO} (\cp{sql.c:6116,6303}); a c\'elula esvaziada n\~ao \'e um
  nulo, \'e uma linha que deixou de fazer parte de $I$.
\end{itemize}
No que diz respeito \`a coluna-alvo de \cp{sql_hist_jev}, a guarda G2
(Sec.~\ref{sec:guardas}) exige $\mathrm{pres}(i,oc)=1$ para toda a linha
viva, pelo que a leitura s\'o corre sobre colunas em que $r_i(oc)$ \'e total.
\end{observacao}

%=======================================================================
\section{N\'ivel 3 — planos, tectos, topo, a c\'elula}
\label{sec:nivel3}

\begin{definicao}[O modelo]
\textbf{Contrato implementado.} O ficheiro de modelo \'e lido por
\cp{mem_le(slot)}, que devolve uma \emph{palavra} ISA de dois \'atomos de 8
bits (\cp{sql.c:1213-1221}); uma zona \'e \cp{ZONA(k)=k<<14}
(\cp{sql.c:208}) com \cp{ZONA_SLOTS=16384} palavras (\cp{sql.c:209}).
Os planos relevantes:
\[
\cp{S_LINHAS}=\cp{ZONA(1)},\quad
\cp{S_DEN}=\cp{ZONA(2)},\quad
\cp{S_ALTO}=\cp{ZONA(3)}
\quad(\cp{sql.c:753,751,834}),
\]
e o tecto de c\'elulas \'e $\cp{CEL\_TECTO}=\cp{S\_ALTO\_N}=16384$
(\cp{sql.c:865-866}). Os bitmaps s\~ao \emph{bit por linha} ou \emph{bit por
c\'elula}:
\[
\cp{S\_VIVO}\subset[512,1024)\ \text{(bit por linha)},\quad
\cp{S\_MATCH}\subset[256,512)\quad(\cp{sql.c:492,495,538})
\]
\[
\cp{S\_PRES}\ \text{(bit por c\'elula)}.
\]
com \cp{LIN\_VIVO}=$512\cdot16=8192$ bits, \cp{LIN\_MATCH}=$4096$ e
\cp{LIN\_TECTO}=\cp{min}=$4096$ (\cp{sql.c:513-516}). O tamp\~ao de nomes
comporta \cp{S_COLNOME_N}=\cp{COL\_MAX}=$1024$ colunas (\cp{sql.c:225,960}).
No \cp{INSERT} recusam-se as linhas acima do tecto
(\cp{sql.c:3649-3654}) e o produto linhas$\times$colunas acima do tecto de
c\'elulas (\cp{sql.c:3659-3664}).
\end{definicao}

\begin{definicao}[Linhas vivas e \emph{topo}]
\label{def:vivo}
\textbf{Contrato implementado.} O \emph{topo} \'e $\cp{nr}=\cp{cat\_nrows()}$
(\cp{sql.c:1504}), o apontador high-water de linhas — n\~ao uma contagem. O
conjunto de linhas vivas \'e
\[
I=\{i\in\Un{nr}:\cp{S\_VIVO}(i)=1\},
\]
e o cardinal $\lvert I\rvert$ s\'o se obt\'em por varredura dos bits
(\cp{sql.c:16354-16357}). Em geral $\lvert I\rvert\ne\nr$: ap\'os um
\cp{DELETE}, $\nr$ fica inalterado e a linha viva deixa de contar.
$\lvert I\rvert=0$ com $\nr>0$ n\~ao \'e erro aqui; a guarda E7 de \emph{tabela
vazia} s\'o dispara para $\nc\le0$ ou $\nr\le0$ (\cp{sql.c:16308}).
\end{definicao}

\begin{definicao}[Valor de c\'elula]
\label{def:celula}
\textbf{Contrato implementado.} Para a c\'elula de endere\c{c}o
$k=i\cdot\nc+j$ (\cp{cel_ix}), o valor lido por \cp{celula_valor}
(\cp{sql.c:7260-7263}) \'e
\[
v(i,j)=w_0{.}\mathrm{total}\;\big|\;(w_1{.}\mathrm{total}\ll 8),
\]
\[
w_0=\cp{mem\_le}(\cp{S\_LINHAS}+k),\qquad
w_1=\cp{mem\_le}(\cp{S\_ALTO}+k),
\]
o inteiro sem sinal de 16 bits em $[0,65535]$. Esta leitura \textbf{n\~ao
consult\~a} \cp{S_PRES} (a presen\c{c}a decide-se \`a parte, guarda G2), e
ignora o segundo \'atomo de cada palavra (o denominador / coeficiente).
O envelope do \cp{INSERT} para \cp{CORPO\_INTEIRO} \'e exactamente
$[0,65535]$ (\cp{sql.c:3665-3690}).
\end{definicao}

%=======================================================================
\section{A leitura: guardas e erros}
\label{sec:guardas}

\textbf{Contrato implementado.} A ordem exacta de guardas de
\cp{sql_hist_jev} (todas em recusa $\Rightarrow$ retorno $0$ e
\cp{out->ok=0}) \'e:
\[
\cp{G7}\xrightarrow{\ \text{E7}\ }\cp{G1}\xrightarrow{\ \text{E1}\ }\cp{G2}\xrightarrow{\ \text{E2}\ }
\cp{G3}\xrightarrow{\ \text{E3}\ }\cp{spec\_decode}\xrightarrow{\ \text{E6}\ }\cp{G4}
\]
\[
\cp{G4}\xrightarrow{\ \text{E4 (inalcan\c{c}avel via \cp{spec})}\ }\cp{G5}\xrightarrow{\ \text{E5}\ }
\cp{F17}\xrightarrow{\ \text{E3}\ }\cp{G6},
\]
seguida da constru\c{c}\~ao da arena, \cp{jev_marginal(X)}, a verifica\c{c}\~ao
de conserva\c{c}\~ao e a sa\'ida. Numericamente:

\begin{center}
\begin{tabular}{@{}l l l@{}}
\toprule
c\'odigo & condi\c{c}\~ao & ponto exacto (\cp{sql.c}) \\
\midrule
E7 & tabela n\~ao existe / coluna n\~ao existe / $\nc\le0$ ou $\nr\le0$ &
16399, 16303, 16308 \\
E1 & corpo da coluna $\cp{oc}\ne\cp{CORPO\_INTEIRO}$ ($=0$) &
16314 \\
E2 & linha viva com $\cp{pres}(i_2,oc)=0$ (\emph{presen\c{c}a incompleta}) &
16323 \\
E3 & valor vivo $v\notin[0,16384)$ (limite do tamp\~ao) &
16338 \\
--  & \cp{spec\_decode} falha (spec malformada) & 16359 \\
E4 & $k\notin[1,16]$ & 16369 \\
E5 & $X<1$ ou $X+2k>16230$ & 16376 \\
E3 & valor vivo $v\ge X$ (F17: limite do suporte) & 16393 \\
E6 & fronteira $b_i\notin(0,X)$ ou n\~ao estritamente crescente & 16402 \\
\bottomrule
\end{tabular}
\end{center}

\begin{observacao}[Contadores]
\textbf{Contrato implementado.} \cp{jev_hist_erros[7]} (\cp{sql.c:16224})
conta E1..E6; \'e \cp{static} e nunca \'e exportado — n\~ao faz parte do
observ\'avel. E7 n\~ao tem contador (\'indice 6 nunca incrementado); E6 usa o
\'indice 5 nas duas ocorr\^encias; E3 usa o \'indice 2 tamb\'em para F17.
\end{observacao}

\begin{observacao}[A ordem \'e contratual]
\textbf{Contrato implementado.} E7 vem antes de G1: uma tabela vazia
($\nr=0$) responde \cp{"E7: tabela vazia"} e nunca G1/E1 — o teste
\cp{test_sql_hist_jev.c} cria uma coluna \cp{TEXTO} com uma linha justamente
para que E1 seja alcan\c{c}avel. G2 vem depois de G1: s\'o se fala de
presen\c{c}a depois de se saber o corpo.
\end{observacao}

%=======================================================================
\section{N\'ivel 1 — realiza\c{c}\~ao, campo e marginal}
\label{sec:nivel1}

Fixam-se, daqui em diante, a tabela $T=(I,J,\mathrm{pres},v)$, a coluna
$\cp{oc}\in J$ e as guardas G1-G6/F17 vencidas.

\begin{definicao}[Realiza\c{c}\~ao]
\label{def:realizacao}
\textbf{Interpreta\c{c}\~ao matem\'atica da implementa\c{c}\~ao.} Gra\c{c}as a
G2, $r_i(\cp{oc})$ \'e total em $I$; gra\c{c}as a G3 e F17, o seu valor vive em
$\Un{X}=\{0,\ldots,X-1\}$. Define-se a \emph{realiza\c{c}\~ao}
\[
\pi\colon I\to\Un{X},\qquad \pi(i)=r_i(\cp{oc}),
\]
e o \emph{campo de contagem}
\[
G\colon\Un{X}\to\N,\qquad G(x)=\lvert\pi^{-1}(x)\rvert,
\]
estendido por $0$ a todo o $V$. Esta \'e exactamente a constru\c{c}\~ao do
\cp{hist_local[16384]} (\cp{sql.c:16338-16352}), que conta com
\cp{celula\_valor}, uma entrada por valor em $[0,16384)$.
\end{definicao}

\begin{definicao}[Parti\c{c}\~ao por fronteiras]
\label{def:particao}
\textbf{Interpreta\c{c}\~ao matem\'atica da implementa\c{c}\~ao.} A \cp{spec}
codifica, por G6, uma sequ\^encia estritamente crescente
\[
B=(b_0,b_1,\ldots,b_k),\qquad b_0=0< b_1<\cdots<b_{k-1}< b_k=X,
\]
e com ela a parti\c{c}\~ao de $\Un{X}$ em $k$ intervalos
\[
q\colon\Un{X}\to\Un{k},\qquad q(x)=j \iff b_j\le x<b_{j+1}.
\]
Como $b_0=0$ e $b_k=X$, $q$ \'e total e cada classe \'e n\~ao vazia.
\end{definicao}

\begin{definicao}[Marginal]
\label{def:marginal}
\textbf{Interpreta\c{c}\~ao matem\'atica da implementa\c{c}\~ao.}
\[
G_q\colon\Un{k}\to\N,\qquad
G_q(j)\;=\;\sum_{x\in q^{-1}(j)}G(x)\;=\;\lvert\{i\in I:q(\pi(i))=j\}\rvert .
\]
\end{definicao}

\begin{lema}[Conserva\c{c}\~ao; duas somas]
\label{lem:conserva}
$\displaystyle\sum_{x\in\Un{X}}G(x)=\lvert I\rvert$
\quad e \quad
$\displaystyle\sum_{j\in\Un{k}}G_q(j)=\lvert I\rvert$.
\begin{proof}
As fibras $\pi^{-1}(x)$ particionam $I$ (aplica\c{c}\~ao total), logo
$\sum_x\lvert\pi^{-1}(x)\rvert=\lvert I\rvert$. A segunda identidade segue de
$\sum_j\sum_{x\in q^{-1}(j)}G(x)=\sum_{x\in\Un{X}}G(x)$, pois $q$ particiona
$\Un{X}$.
\end{proof}
\end{lema}

\begin{proposicao}[Correspond\^encia com a arena]
\label{prop:corr}
\textbf{Contrato implementado.} Sob as mesmas guardas, com $\cp{spec}$ a dar
$(X,k,B)$, o c\'odigo monta na arena (Sec.~\ref{sec:arena}) e chama
\cp{jev_marginal(X)}. Dado $D_1=X+k+10$, o resultado da primitiva satisfaz
\[
\cp{a}[D_1+j]=G_q(j)\quad(j\in\Un{k}),
\]
e a soma de conserva\c{c}\~ao \'e verificada por
$\sum_{j<k}\cp{a}[D_1+j]=\lvert I\rvert$ (\cp{sql.c:16461-16470}), onde
$\lvert I\rvert$ veio de \cp{a[X]}=\cp{*pi1}.
\begin{proof}
Em \cp{jev_marginal_raw} (\cp{jev.c:97-141}) com
$(c_0,c_1,c_2)=(k,1,1)$, para cada $x\in\Un{X}$:
$y_0=\cp{proj}(b_0..b_k,c_0,x)=q(x)\in\Un{k}$, e $y_1=y_2=0$ pois ambos os
blocos s\~ao $\{0,X\}$. Logo $\cp{a}[D_1+y_0]\mathrel{+}=G(x)$; como $q(\cdot)$
avalia \`a \'antiga de cada $x$ que contribui, o acumulador na posi\c{c}\~ao $j$
recebe precisamente $\sum_{x\in q^{-1}(j)}G(x)$. A conjunta
$J[\cdot]$ com \'indice $\cp{idx}=y_2+c_2(y_1+c_1y_0)=y_0$ recebe o mesmo, e por
isso $D_1[0..k)=J[0..k)$ (a segunda passagem, o od\'ometro,
\cp{jev.c:127-137}, n\~ao os altera — ver Sec.~\ref{sec:arena}). No bloco
\cp{OUT}: $\cp{a[OUT]}=\sum_x G(x)=\lvert I\rvert$ e
$\cp{a[OUT+1]}=\cp{a[X]}=\lvert I\rvert$.
\end{proof}
\end{proposicao}

\begin{observacao}[A verifica\c{c}\~ao l\^e a pr\'opria sa\'ida]
\textbf{Interpreta\c{c}\~ao matem\'atica da implementa\c{c}\~ao.} A guarda de
conserva\c{c}\~ao n\~ao recalcula $G_q$ por outra via: soma os valores de
\cp{D1} que a primitiva escreveu e compara-os com $\lvert I\rvert$. \'E uma
auto-consist\^encia estrutural apoiada no Lema~\ref{lem:conserva} — e, na
hist\'oria do ficheiro, foi ela que apanhou G nulo (F9) e $D_1$ deslocado
(F10), como o pr\'oprio coment\'ario de \cp{sql.c:16455-16460} regista.
\end{observacao}

%=======================================================================
\section{A \cp{spec} base-36}
\label{sec:spec}

\begin{definicao}[Decodificador]
\textbf{Contrato implementado.} \cp{spec_decode} (\cp{sql.c:16233}) l\^e
$s=\cp{spec}$ sob o alfabeto base-36 ($0\ldots9,a\ldots z$; mai\'usculas s\~ao
d\'igito inv\'alido, \cp{base36\_digit} devolve $-1$). O comprimento \'e da
forma
\[
\ell = 4+3(k-1), \qquad k=\frac{\ell-4}{3}+1\in[1,16],
\]
e os blocos s\~ao: \'indices $0..2$ o valor $X=d_0\cdot1296+d_1\cdot36+d_2$
($X\ge1$); \'indice $3$ uma redund\^ancia: o d\'igito tem de ser igual a $k$
(\cp{sql.c:16246-16247}); cada triplo $3j..3j+2$, $j=1..k-1$, decifra a
fronteira $b_j$. O decodificador n\~ao valida as fronteiras (valores at\'e
$46655$): essa valida\c{c}\~ao \'e G6, logo ap\'os ($0<b_j<X$ e estritamente
crescentes). A recusa do decodificador \'e \cp{"E6: spec invalida"}
(\cp{sql.c:16362-16365}).
\end{definicao}

\begin{definicao}[O objecto codificado]
\textbf{Interpreta\c{c}\~ao matem\'atica da implementa\c{c}\~ao.} A \cp{spec}
\emph{sob guardas vencidas} codifica a tripla
\[
(X,\ k,\ B=(b_1,\ldots,b_{k-1}))
\]
da Sec.\ref{sec:nivel1}. A redund\^ancia do 4.\'o d\'igito torna o
comprimento observ\'avel pelo pr\'oprio c\'odigo: um \cp{spec} com
$\ell=4+3(k'-1)$ e d\'igito $d_3\ne k'$ \'e rejeitado mesmo que $\ell$ seja
formalmente v\'alido.
\end{definicao}

Nota: o exemplo can\'onico dos testes \'e \cp{"00a3002007"}, que codifica
$X=10$, $k=3$, $B=(b_1,b_2)=(2,7)$ — com $b_0=0$ e $b_3=10$ impl\'icitos.

%=======================================================================
\section{A arena e a constante $16230$}
\label{sec:arena}

\begin{definicao}[Arena]
\textbf{Contrato implementado.} \cp{jev_arena()} devolve o \'unico
\cp{static unsigned char arena[65536]} (\cp{jev.c:83,319-321}), lido como
$\cp{JEV\_ARENA\_BYTES}/4=2^{14}=16384$ inteiros \cp{int32}
$\cp{a}[0..16383]$ (\cp{jev_api.h:36}). N\~ao \'e o ficheiro de modelo: \'e
outro buffer, em mem\'oria nativa, de 16\,384 palavras de 32 bits — coincidindo
em n\'umero com \cp{ZONA\_SLOTS} sem ser o mesmo objecto.
\end{definicao}

\begin{definicao}[Layout da marginal com $(k,1,1)$]
\label{def:layout}
\textbf{Contrato implementado.} O bloco que \cp{sql_hist_jev} monta
(\cp{sql.c:16417-16442}) e que \cp{jev_marginal_raw} interpreta
(\cp{jev.c:97-141}) — numerando as c\'elulas:
\begin{center}
\begin{tabular}{@{}l l l@{}}
\toprule
\'indice & conte\'udo & origem \\
\midrule
$[0,X)$ & $G(x)=\cp{a}[x]$ & \cp{hist\_local} \cp{sql.c:16422} \\
$X$ & $\lvert I\rvert$ & \cp{*pi1} \cp{sql.c:16423} \\
$X+1$ & $0$ ($n$; s\'o \cp{marginal\_d} o l\^e) & fixo \cp{sql.c:16424} \\
$X+2..X+4$ & $c_0=k,\ c_1=1,\ c_2=1$ & F10 \cp{sql.c:16425} \\
$X+5$ & $b_0=0$ & fixo \\
$X+5+j\ (1\le j<k)$ & $b_j$ & \cp{bspec[j-1]} \\
$X+5+k$ & $b_k=X$ & fixo \\
$X+5+k+1..X+5+k+4$ & $0,X,0,X$ (blocos 1 e 2) & fixo \cp{sql.c:16429-16430} \\
$D_1=X+k+10$ & marginais directas, linha $p$ em $16p$ & \cp{jev_marginal} \\
$J=D_1+48$ & conjunta, $\cp{PROD}=k$ c\'elulas & \cp{jev_marginal} \\
$D_2=J+\cp{PROD}$ & marginais da conjunta & \cp{jev_marginal} \\
$\cp{OUT}=D_2+48$ & $\cp{OUT[0]}$=total, $\cp{OUT[1]}=\lvert I\rvert$ &
\cp{jev_marginal} \\
\bottomrule
\end{tabular}
\end{center}
A express\~ao de $D_1$ reproduz a da primitiva (n\~ao \'e uma constante
escolhida): $\cp{P}=X+2$, $\cp{B}=\cp{P}+3$, $\cp{S}=c_0+c_1+c_2+3=k+5$, e
$D_1=\cp{B}+\cp{S}=X+k+10$ (\cp{sql.c:16437-16440}; idem \cp{jev.c:99-103}).
\end{definicao}

\begin{proposicao}[A constante $16230$]
\label{prop:16230}
\textbf{Contrato implementado.} A \'ultima c\'elula que a marginal escreve \'e
$\cp{OUT}+47$, e a condi\c{c}\~ao de caber na arena \'e
\[
\cp{OUT}+47 = X+2k+106+47 = X+2k+153 \le 16383
\iff X+2k\le 16230 ,
\]
que \'e exactamente a guarda G5/E5 (\cp{sql.c:16376}). Deriva\c{c}\~ao:
$J=X+k+58$, $D_2=J+\cp{PROD}=X+2k+58$ (com $\cp{PROD}=c_0c_1c_2=k$),
$\cp{OUT}=D_2+48=X+2k+106$. O resto $16384-16230=154$ \'e o \'onus fixo
$106+48$ acima da fronteira $X+2k$.
\end{proposicao}

\begin{observacao}[Restaura\c{c}\~ao da arena]
\textbf{Contrato implementado.} \cp{sql_hist_jev} grava os $65536$ bytes antes
de tocar na arena (\cp{sql.c:16417-16418}, \cp{memcpy} para \cp{guardado}) e
restaura a c\'opia em todos os caminhos que passam pela constru\c{c}\~ao:
falha de \cp{jev_marginal} (\cp{16448}), conserva\c{c}\~ao violada
(\cp{16464}) e sucesso (\cp{16481}); os caminhos anteriores a 16417 nunca
tocam na arena. O resultado \'e a transpar\^encia da arena. O coment\'ario do
cabe\c{c}alho \cp{sql.c:16223-16224} diz que o snapshot \'e \emph{por conta do
chamador} — o c\'odigo \'e mais forte do que o coment\'ario: faz o pr\'oprio.
\textbf{(N\~ao estabelecido pela implementa\c{c}\~ao atual:} se a inten\c{c}\~ao
era deixar o snapshot ao chamador, o c\'odigo atual n\~ao o deixa.)
\end{observacao}

%=======================================================================
\section{As leituras da porta C}
\label{sec:porta}

\textbf{Contrato implementado.} A porta \cp{jev_api.h} exp\~oe as cinco
primitivas e a projec\c{c}\~ao; todas com retorno $1$ ($\cp{JEV\_OK}$) ou $0$
($\cp{JEV\_ERR\_CLASSE}$), excepto \cp{jev_proj}, que responde $-1$ como
resposta leg\'itima (\cp{jev_api.h:22-31,74-77}; wrappers \cp{jev.c:325-351}):

\begin{center}
\begin{tabular}{@{}>{\raggedright\arraybackslash}p{3.8cm} >{\raggedright\arraybackslash}p{5.4cm} >{\raggedright\arraybackslash}p{5.2cm}@{}}
\toprule
primitiva & interpreta\c{c}\~ao & layout (N\'ivel 3) \\
\midrule
\cp{jev_proj}(b,c,x) & classe de $x$ em $c+1$ fronteiras $b_0..b_c$; $-1$ se fora & varredura de compara\c{c}\~oes, \cp{jev.c:85-95} \\
\cp{jev_marginal}(X) & marginais da conjunta igualam as directas ($\sum_{j}D_2=\sum_j D_1$); $\cp{OUT[0]}$=total, $\cp{OUT[1]}=\lvert I\rvert$ & $c_0c_1c_2$, od\'ometro, \cp{jev.c:97-141} \\
\cp{jev_escore}(X) & funcional $N=\sum_{i}s_i\,G_q(s_i)$ como par racional ex\'acto $(N,\lvert I\rvert)$, sem divis\~ao & $m=\cp{a[X+1]}$, pesos em $\cp{W}=\cp{B}+m+1$, $\cp{OUT[0..2]}=N,\lvert I\rvert,m$; \cp{jev.c:156-171} \\
\cp{jev_duas}(X) & Noul e Score na mesma varredura sobre o mesmo $G$ & \cp{jev.c:191-...} \\
\cp{jev_marginal\_d}(X) & projec\c{c}\~oes como coordenadas de $\mathcal{X}_d$ & \cp{jev.c:238-...} \\
\bottomrule
\end{tabular}
\end{center}

\begin{observacao}[O caminho SQL usa duas]
\textbf{Contrato implementado.} O caminho \cp{sql_hist_jev} chama
\cp{jev_arena} e \cp{jev_marginal} apenas. As outras leituras existem na
porta para o produto, mas este ficheiro n\~ao as invoca; \emph{em particular,
\cp{jev_escore} n\~ao \'e chamado}, logo o funcional $L(p)$ n\~ao interv\'em no
resultado SQL (ver Sec.~\ref{sec:naoest}).
\end{observacao}

%=======================================================================
\section{Verifica\c{c}\~ao}
\label{sec:verif}

\begin{itemize}
\item \medido{\cp{tests/test\_sql\_hist\_jev.c} exercita o caminho completo
  \cp{tabela -> linhas vivas -> hist\_local -> arena G -> jev\_marginal ->
  D1 -> SqlOut}. O caso nominal T1: $X=10$, $B=\{0,2,7,10\}$, valores $0..9$
  ($\lvert I\rvert=10$) devolve $G_q=(2,5,3)$ e soma $10$. T2 (fronteira
  $v=X-1$), T3/T4 (F17: $v=X$ e $v>X$ recusados com E3), T5 (E3, limite do
  tamp\~ao: $v=20000$), T6 (E1: corpo \cp{TEXTO}), T7 (E7).
  Cada chamada sela a arena com um padr\~ao alternante ($0\mathsf{x}A5\oplus i$)
  e compara-a byte a byte depois — prova de restaura\c{c}\~ao.}
\item \medido{\cp{FASE\_4\_9\_R3\_RESULTADO.md}: 11 casos sobre o mesmo
  harness. PASS 7/11 (E7, E4/$k>16$, E5/$X+2k>16230$, E6/n\~ao crescente,
  B/$X{=}10\ B{=}\{0,3,8,10\}$, C/$k{=}1\ X{=}1$, setup); 2 FAIL eram
  anomalias dos checks do harness (A comparava \cp{gq} com \cp{'0'} apesar de
  $G_0=2$; D usava $b_1=0$, proibido por G6); 3 NI (E1, E2, E3 sem tabela
  populada). Conclus\~ao do relat\'orio: os checks falhavam, n\~ao o
  \cp{sql.c}.}
\item \medido{Re-deriva\c{c}\~ao independente desta auditoria: um harness
  avulso, ligado s\'o a \cp{banco/jev.c}, monta o layout da
  Def.~\ref{def:layout} e confirma $D_1=X+k+10=23$, $G_q=(4,12,9)$ contra uma
  soma intervalar independente para $X=10$, $\lvert I\rvert=25$,
  $\cp{OUT[0]}=\cp{OUT[1]}=25$, e \cp{jev_marginal}=\cp{JEV_OK}.}
\end{itemize}

%=======================================================================
\section{N\'ivel 2 --- opera\c{c}\~oes relacionais implementadas}
\label{sec:relacional}

Esta sec\c{c}\~ao modela as opera\c{c}\~oes que \cp{sql.c} realmente executa,
como objectos e opera\c{c}\~oes matem\'aticas sobre o modelo da
Def.~\ref{def:tabela}. S\'o entra o que o c\'odigo faz; a matem\'atica nasce do
c\'odigo e n\~ao de um tratado de \'algebra relacional. Para cada opera\c{c}\~ao
diz-se a defini\c{c}\~ao (N\'ivel 1), o que a implementa\c{c}\~ao sustenta
(N\'ivel 3) e o estado de correspond\^encia. A leitura JEV n\~ao \'e alterada:
ela \'e o \'ultimo elo da cadeia e permanece como nas
Secs.~\ref{sec:nivel1}--\ref{sec:porta}.

\subsection{Os quatro estados}
Os tr\^es carimbos das Secs.~\ref{sec:guardas} (marcas de
\textbf{Contrato implementado.}, \textbf{Interpreta\c{c}\~ao matem\'atica da
implementa\c{c}\~ao.} e \textbf{N\~ao estabelecido pela implementa\c{c}\~ao
atual.}) continuam a valer. Para auditar
opera\c{c}\~ao a opera\c{c}\~ao, o primeiro deles refina-se em tr\^es estados,
que somam quatro com o terceiro carimbo:

\begin{center}
\begin{tabular}{@{}p{6.2cm} p{9.0cm}@{}}
\toprule
estado & leitura \\
\midrule
\sti & o \textbf{Contrato implementado.} com a constru\c{c}\~ao matem\'atica
escrita e ancorada no c\'odigo; \\
\stipa & o \textbf{Contrato implementado.} vale, mas o canal de sa\'ida ou o
protocolo n\~ao determinam sozinho a matem\'atica --- o conflito fica
declarado, n\~ao resolvido por suposi\c{c}\~ao; \\
\stco & o comportamento \'e compat\'ivel com a leitura, mas este paper n\~ao o
demonstra; \\
\stni & o que a implementa\c{c}\~ao atual n\~ao faz (o terceiro carimbo, sem
altera\c{c}\~ao). \\
\bottomrule
\end{tabular}
\end{center}

\subsection{A tabela como multiconjunto}
\label{sec:mult}
\begin{definicao}[Multiconjunto de linhas]
\label{def:multiconjunto}
\textbf{Interpreta\c{c}\~ao matem\'atica da implementa\c{c}\~ao.} A tabela $T$
da Def.~\ref{def:tabela} lida como objecto relacional \'e o \emph{multiconjunto}
das suas linhas vivas,
\[
R=\big(\,r_i\;\big)_{i\in I},
\]
indexado pelo conjunto $I$ (Def.~\ref{def:vivo}). Duas linhas com o mesmo
conte\'udo coexistem: a identidade \'e o \'indice $i$, nunca o valor. Nada no
c\'odigo funde linhas iguais automaticamente --- dois \cp{INSERT} do mesmo
\cp{VALUES} d\~ao duas linhas vivas ---, e a fus\~ao por conte\'udo \'e
expl\'icita e s\'o acontece na sa\'ida de \cp{DISTINCT} ou \cp{UNION}
(Secs.~\ref{sec:proj} e \ref{sec:uni}). A multiplicidade \'e o que faz do
campo $G$ uma contagem: na coluna $\cp{oc}$, $G(x)=\#\{i\in I:r_i(\cp{oc})=x\}$
(Def.~\ref{def:realizacao}).
\end{definicao}

\begin{observacao}[Corpo e envelope]
\textbf{Interpreta\c{c}\~ao matem\'atica da implementa\c{c}\~ao.} Cada coluna
$j$ declara um \emph{corpo} $D_j$ --- um dos oito do \cp{CREATE TABLE}
(\cp{Inteiro, Racional, \'Aureo, M\'orfico, Cristalino, Booleano, Texto, Data};
\cp{sql.c:3066-3546}) --- e um \emph{envelope}: a g\'ama declarada que aceita. O
valor de c\'elula vive em $V_j=D_j$: a palavra de dois andares da
Def.~\ref{def:celula} (num racional, o primeiro \'atomo \'e o numerador
assinado e o segundo o denominador). A \emph{aus\^encia} \'e o dual (bit
\cp{S_PRES} a $0$, Obs. da Sec.~\ref{sec:tabela}); a leitura de uma c\'elula
ausente devolve o neutro do corpo, mas esse neutro n\~ao \'e um dado --- o que
faz o valor \'e o bit de presen\c{c}a. O envelope recusa em voz alta o que n\~ao
cabe, e a recusa n\~ao altera o estado.
\end{observacao}

\subsection{Selec\c{c}\~ao: \cp{WHERE}}
\label{sec:where}
A condi\c{c}\~ao $\varphi$ \'e lida por \cp{le_expr} (\cp{sql.c:4798}):
\'atomos ``$L\;op\;R$'' com $op\in\{=,\ne,<,\le,>,\ge\}$, compostos por
\cp{AND}, \cp{OR} e \cp{NOT} (este empurrado para as folhas por De Morgan,
\cp{nega_arvore}, \cp{sql.c:4669-4681}); \cp{x IN (v1, v2, ...)} \'e a
disjun\c{c}\~ao das igualdades
($\varphi=(x{=}v_1)\lor\cdots\lor(x{=}v_m)$), \cp{sql.c:4704-4749}); e
\cp{x IS [NOT] NULL} \'e uma cl\'ausula \`a parte, lida do bit \cp{S_PRES} e
n\~ao compilada para a ISA (\cp{sql.c:8232-8262, 8988-8994}).

Seja $C(\varphi)\subseteq J$ o conjunto das colunas que $\varphi$ \emph{cita}
(o \cp{citadas_where} do motor) e $\tilde r_i$ a linha $r_i$ estendida pelo
neutro onde est\'a ausente. A regra efectiva \'e determinada pela
implementa\c{c}\~ao, e \'e \emph{bivalente}:
\begin{definicao}[Selec\c{c}\~ao]
\label{def:selecao}
\textbf{Interpreta\c{c}\~ao matem\'atica da implementa\c{c}\~ao.} Para cada
linha viva $i\in I$,
\[
i\in I(\varphi)
\;\;\Longleftrightarrow\;\;
\Big(\ \forall j\in C(\varphi):\ \mathrm{pres}(i,j)=1\ \Big)
\;\wedge\;
\varphi\big(\tilde r_i\big)=\mathbf{t},
\]
e a selec\c{c}\~ao \'e
\[
\sigma_\varphi(R)=\big(\,r_i\;\big)_{i\in I(\varphi)} .
\]
O molde emitido por linha avalia $\varphi$ sobre $\tilde r_i$ (as oito
rela\c{c}\~oes de \cp{contrato.h} e a contra\c{c}\~ao \cp{fecha_cond},
\cp{sql.c:4633-4653}, que decide na compila\c{c}\~ao o que n\~ao depende da
linha); depois, uma passagem final apaga toda a linha que cita uma coluna
ausente (\cp{sql.c:8933-8941}). O resultado n\~ao depende do que esteja escrito
numa c\'elula ausente: essa linha j\'a n\~ao casa.
\end{definicao}

Consequ\^encias determinadas pela mesma leitura:
\begin{itemize}
\item \textbf{A compara\c{c}\~ao com a aus\^encia n\~ao casa.} \emph{Comparar
  \'e pedir um valor, e a aus\^encia n\~ao o tem} (\cp{sql.c:8792-8799}). Uma
  condi\c{c}\~ao que cite uma c\'elula ausente exclui a linha, em qualquer
  polaridade: $\lnot(A\lor B)$ com $B$ a citar a aus\^encia exclui tanto quanto
  $A\lor B$. \stni{} l\'ogica tern\'aria: n\~ao existe \texttt{UNKNOWN}, e o
  \texttt{NULL} do SQL n\~ao \'e aqui um valor de verdade.
\item \textbf{\cp{IS [NOT] NULL} s\~ao os \'unicos testes do dual.} Mant\^em ou
  deitam fora as linhas conforme o bit \cp{S_PRES} da coluna (\cp{nul_nega},
  \cp{sql.c:8988-8994}); a coluna testada n\~ao entra em $C(\varphi)$.
\item \textbf{\cp{IN} e \cp{NOT IN} sobre subconsulta.} A linha cuja chave \'e
  ausente \'e retirada \emph{antes} da perten\c{c}a, nos dois sentidos --- ``a
  aus\^encia n\~ao pertence nem deixa de pertencer'' (\cp{sql.c:8862-8869}); a
  perten\c{c}a decide-se sobre o corpo da subconsulta, e o \cp{NOT IN} \'e o
  complemento sobre essas linhas (\cp{sql.c:8903-8909}).
\item \textbf{\cp{EXISTS} \'e um quantificador.} Decide-se uma vez sobre a
  tabela da subconsulta (existe linha viva l\'a?) e \'e por isso uma constante
  para a consulta; \cp{NOT EXISTS} \'e o complemento --- a mesma De Morgan
  sobre um conjunto (\cp{sql.c:8280-8331, 8982-8986}).
\item \textbf{Poda por alcance.} Com um s\'o \'atomo, sem agrega\c{c}\~ao nem
  agrupamento, a guarda \cp{atom_possivel} (\cp{sql.c:8571-8576}) esvazia o
  campo sem correr molde quando a gama que as colunas citadas deixaram escrita
  torna o \'atomo imposs\'ivel: \'e uma optimiza\c{c}\~ao que preserva a
  resposta, n\~ao uma sem\^antica nova.
\end{itemize}

\textbf{Status: \sti} A leitura de \cp{le_expr} e do percurso de
\cp{sql.c:8710-8841, 8933-8941} determina uma \'unica sem\^antica: a bivalente
da Def.~\ref{def:selecao}. N\~ao se inventou l\'ogica tern\'aria.

\subsection{Projec\c{c}\~ao e \cp{DISTINCT}}
\label{sec:proj}
\begin{definicao}[Projec\c{c}\~ao]
\textbf{Interpreta\c{c}\~ao matem\'atica da implementa\c{c}\~ao.} A lista do
\cp{SELECT} \'e um mapa $p\colon\Un{m}\to J$ (ou express\~oes). Dada a
selec\c{c}\~ao $\sigma_\varphi(R)$, a projec\c{c}\~ao devolve a fam\'ilia
\[
\pi_{p}\big(\sigma_\varphi(R)\big)=\big(\ \rho_{i}\ \big)_{i\in I(\varphi)},
\qquad
\rho_i(k)=\beta_k(r_i),
\]
onde $\beta_k$ \'e o valor da coluna $p(k)$ ou o polin\'omio (tensor) da
express\~ao sobre as colunas citadas --- at\'e \cp{NCOL}$=6$ colunas por
express\~ao, e o excesso \'e recusado (\cp{sql.c:4170, 4540-4543}). Se alguma
coluna citada pela express\~ao est\'a ausente em $i$, a c\'elula $\rho_i(k)$ \'e
emitida como ausente --- a mesma regra do \cp{CHECK} e do \cp{WHERE}, ``se o
resultado d\'a com uma coluna ausente, o resultado \'e ausente''
(\cp{sql.c:13048-13050}).
\end{definicao}

\textbf{Contrato implementado.} \cp{SELECT <constante>} (sem \cp{FROM}) \'e o
caso degenerado da projec\c{c}\~ao: a express\~ao sem vari\'aveis, que devolve
uma linha com o valor, e o caminho por que um cliente testa a liga\c{c}\~ao
(\cp{sql.c:15899-15927}).

\begin{definicao}[Distinct de uma coluna]
\textbf{Interpreta\c{c}\~ao matem\'atica da implementa\c{c}\~ao.} O
\cp{DISTINCT} devolve a \emph{imagem da fibra}:
\[
\operatorname{DISTINCT}_d\big(\sigma_\varphi(R)\big)
=\big\{\, r_i(d) : i\in I(\varphi) \,\big\},
\]
com um representante can\'onico por valor --- ``a folha 1 de cada fibra''
(\cp{sql.c:12857-12861}) --- e a aus\^encia como \emph{um} valor distinto, que
n\~ao se junta ao zero escrito. \textbf{Contrato implementado:} pede uma
coluna de cada vez --- ``a linha inteira n\~ao \'e index\'avel pela \'arvore'' e
\'e recusada com o n\'umero (\cp{proj_n != 1}, \cp{sql.c:12846-12854}). A
deduplica\c{c}\~ao \'e a inser\c{c}\~ao na \'arvore do corpo e \cp{ord_tira_repetidos}
(\cp{sql.c:12885-12898}); se a \'arvore n\~ao couber, recusa (nunca devolve
repetidos).
\end{definicao}

\textbf{Status das duas: \sti} com as particularidades declaradas (uma coluna
por \cp{DISTINCT}; a aus\^encia \'e uma classe).

\subsection{Agrega\c{c}\~ao e contagem}
\label{sec:agreg}
\begin{itemize}
\item \textbf{\cp{COUNT(*)}} \'e a cardinalidade $|I(\varphi)|$, \emph{lida} e
  n\~ao reconta: soma-se o campo do match (o popcount), ``a leitura da que j\'a
  ficou escrita'' (\cp{sql.c:15929-15935}); numa tabela vazia \'e zerado
  explicitamente (\cp{sql.c:8664-8679}). Zero aqui \'e o valor correcto: contar
  nada \'e zero, n\~ao uma aus\^encia.
\item \textbf{\cp{COUNT(DISTINCT d)}} \'e o n\'umero de classes do quociente por
  $d$ sobre $I(\varphi)$ --- o mesmo objecto do \cp{GROUP BY} lido em
  quantidade --- e a aus\^encia \'e uma classe, uma s\'o (\cp{sql.c:9000-9044}).
\item \textbf{\cp{SUM}, \cp{MAX}, \cp{MIN}, \cp{AVG}} correm na mesma passagem
  com um acumulador por agrega\c{c}\~ao (\cp{sql.c:12713-12745}) e somam o
  \emph{corpo}, n\~ao o suporte: c\'elulas ausentes n\~ao alimentam
  (\cp{sql.c:12739}). \cp{AVG} devolve o par racional reduzido (a classe racional
  do corpo, \cp{sql.c:12708}). ``A soma de nada \'e \emph{ausente}, n\~ao zero.''
\item \textbf{A agrega\c{c}\~ao sem \cp{GROUP BY} \'e o quociente pelo mapa
  constante:} uma fibra s\'o, a tabela inteira (\cp{sql.c:9766-9777}).
\end{itemize}

\textbf{Status: \sti}

\subsection{\cp{GROUP BY} e \cp{HAVING}}
\label{sec:grp}
A chave de agrupamento é dada por uma ou duas colunas. Para uma coluna,
$k(i)=r_i(c)$; para duas, $g(i)=(k_1(i),k_2(i))$ com
$k_1(i)=r_i(c_1), k_2(i)=r_i(c_2)$. A construção é por fibras (Sec.~\ref{sec:fibras}); a mesma árvore partilha a estrutura de indexação (\cp{sql.c:12612-12656}).

\begin{definicao}[Agrupamento]
\textbf{Interpreta\c{c}\~ao matem\'atica da implementa\c{c}\~ao.}
Seja $k$ a função de chave sobre $I(\varphi)\cap I'$. Defina
\[
\operatorname{GROUP}(I(\varphi),k)
=
\{\operatorname{Fib}(k,v): v\in\operatorname{im}(k)\}.
\]
Para chave composta,
\[
\operatorname{Fib}(g,(v_1,v_2))
=
\{i\in I(\varphi)\cap I': g(i)=(v_1,v_2)\}
=
\operatorname{Fib}(k_1,v_1)\cap\operatorname{Fib}(k_2,v_2),
\]
e
\[
\operatorname{GROUP}(I(\varphi),g)
=
\{\operatorname{Fib}(g,(v_1,v_2)): (v_1,v_2)\in\operatorname{im}(g)\}.
\]
A saída tem uma linha por fibra com a chave. A coluna \cp{count} é
\emph{sempre} emitida com a cardinalidade da fibra
\[
\operatorname{COUNT}(*)=|\operatorname{Fib}(k,v)|\quad\text{(resp. }|\operatorname{Fib}(g,(v_1,v_2))|\text{)}\quad (\cp{sql.c:12703}).
\]
As instâncias cuja chave está ausente ($I_\bot\cap I(\varphi)$) formam a sua
própria classe no fim --- ``o dual é um grupo, e não o grupo do zero''
(\cp{sql.c:12628-12633}) --- e não são elementos de nenhuma $\operatorname{Fib}(k,v)$ ordinária.
As fibras ordinárias enumeram-se por ordem da árvore (\cp{sql.c:12612-12656}).
Conservação: $\sum_{F\in\operatorname{GROUP}(I(\varphi),k)} |F| + |I_\bot\cap I(\varphi)| = |I(\varphi)|$
(\cp{sql.c:12745-12748}).
\end{definicao}

\textbf{Contrato implementado.} O \cp{HAVING} s\'o aceita uma condi\c{c}\~ao
sobre o \cp{count}: $|R_k|>n$, $|R_k|<n$ ou $|R_k|=n$, filtrada depois das
fibras (\cp{sql.c:8393-8429, 12749-12753}). \stni{} \cp{HAVING} geral (sobre
somas, m\'edias ou chaves).

\textbf{Status: \sti}

\subsection{Ordem e a fatia: \cp{ORDER BY}, \cp{LIMIT}, \cp{OFFSET}}
\label{sec:ord}
\begin{definicao}[Ordem]
\textbf{Interpreta\c{c}\~ao matem\'atica da implementa\c{c}\~ao.} O
\cp{ORDER BY} \'e a permuta\c{c}\~ao $\sigma$ de $I(\varphi)$ que ordena as
linhas pela \emph{r\'egua} da primeira coluna, isto \'e, pela chave
$\kappa_1(i)$ que \cp{col_regua_chave}($i$, \cp{c1}, \cp{ncols}) devolve
(\cp{sql.c:9685, 9666-9676}),
ascendente ou descendente conforme a direc\c{c}\~ao da descida da \'arvore
(\cp{ord_desc}); dentro de cada fibra da primeira, corre-se a mesma descida com
a segunda coluna (\cp{sql.c:9703-9752}). A r\'egua \'e a do corpo --- numa
coluna quad\'ratica, a norma $q(a,b)=a^2+B\,ab+C\,b^2$ de assinatura
$\Delta=B^2-4C$, que o motor \emph{declara} numa nota a cada consulta
(\cp{sql.c:9666-9676}) --- n\~ao a ordem dos bits da c\'elula. Se a \'arvore de
ordena\c{c}\~ao n\~ao couber, a consulta \'e recusada (ordenar metade seria
devolver uma ordem que ningu\'em pediu, \cp{sql.c:9612-9615, 9686-9693}).
\end{definicao}

\textbf{Contrato implementado.} As particularidades:
\begin{itemize}
\item \textbf{Texto recusado.} \cp{ORDER BY} sobre coluna de texto \'e
  recusado com a raz\~ao: a c\'elula guarda o \emph{endere\c{c}o} no pool, que
  preserva a igualdade e n\~ao a ordem; ordenar pelo prefixo daria uma ordem
  parcial apresentada como total (\cp{sql.c:9620-9647}).
\item \textbf{Ausentes ao fim, em bloco.} A linha cuja primeira chave \'e
  ausente n\~ao tem com que comparar; vai ao fim, em bloco --- a conven\c{c}\~ao
  do \cp{NULLS LAST}: ``o dual n\~ao \'e pequeno, \'e de outro n\'ivel''
  (\cp{sql.c:9648-9653, 9700}).
\end{itemize}

\begin{definicao}[A fatia]
\textbf{Contrato implementado.} Sobre a sequ\^encia ordenada aplica-se a janela:
o \cp{OFFSET} salta primeiro e o \cp{LIMIT} conta depois --- ``s\~ao os dois
extremos da mesma faixa, e trocar a ordem deles daria outra fatia''
(\cp{sql.c:12906-12910}). Sintaticamente aceitam-se \cp{LIMIT n [OFFSET m]} ou
\cp{OFFSET m} (\cp{sql.c:8504-8531}); o resto \'e recusado. A porta C limita a
materializa\c{c}\~ao: no m\'aximo \cp{SQL\_OUT\_MAX\_ROWS} $=64$ linhas, 64
colunas e 64 bytes por c\'elula (\cp{sql_api.h:82-84}); o excesso n\~ao cabe no
\cp{SqlOut} e corta-se em sil\^encio (\cp{sql.c:12913}).
\end{definicao}

\textbf{Status: \sti}

\subsection{Jun\c{c}\~ao}
\label{sec:join}
Duas tabelas $R=(r_i)_{i\in I_R}$ e $S=(s_t)_{t\in I_S}$ com conjuntos de
colunas disjuntos $J_R$, $J_S$. O \cp{ON} da jun\c{c}\~ao \'e \emph{um}
predicado de uma coluna de cada lado, $\mathrm{x}.\mathrm{a}\;\theta\;\mathrm{y}.\mathrm{b}$
com $\theta\in\{=,<,\le,>,\ge\}$ (sem constantes, sem \cp{<>}, sem \cp{AND}; os
lados trocados invertem o operador, \cp{sql.c:7710-7736}).

\begin{definicao}[Concatena\c{c}\~ao e $\theta$-jun\c{c}\~ao]
\textbf{Interpreta\c{c}\~ao matem\'atica da implementa\c{c}\~ao.} Com colunas
disjuntas, a concatena\c{c}\~ao $r_i\oplus s_t\colon J_R\uplus J_S\dashrightarrow V$
\'e a linha dada por $r_i$ em $J_R$ e por $s_t$ em $J_S$. Definam-se
$I_R'=\{i\in I_R:\mathrm{pres}_R(i,a)=1\}$,
$I_S'=\{i\in I_S:\mathrm{pres}_S(i,b)=1\}$ e as fibras
$\operatorname{Fib}(k_B,v)=\{t\in I_S': k_B(t)=v\}$ (Sec.~\ref{sec:fibras}).

A $\theta$-jun\c{c}\~ao interna é
\[
R\Join^{\theta}_{a,b}\, S
=\big\{\, r_i\oplus s_t :
          i\in I_R',\ t\in I_S',\ k_R(i)\,\theta\, k_B(t) \,\big\},
\]
em que, para igualdade,
\[
\boxed{
R\Join^{=}_{a,b}\, S
=
\bigcup_{i\in I_R'}
\{r_i\oplus s_t : t\in \operatorname{Fib}(k_B,k_R(i))\}
=
\bigcup_{i\in I_R'}
\left(\{r_i\}\times
\operatorname{Fib}(k_B,k_R(i))\right)
}
\]
(vista no conjunto de pares de linhas), e, para desigualdades,
\[
R\Join^{\theta}_{a,b}\, S
=
\bigcup_{i\in I_R'}
\left(
\{r_i\}\times
\bigcup_{\substack{v\in\operatorname{im}(k_B)\\
v\theta k_R(i)}}
\operatorname{Fib}(k_B,v)
\right).
\]
O produto tem multiplicidade por pares: se dois índices $t_1\neq t_2$
coexistem em $\operatorname{Fib}(k_B,v)$, ambos casam com cada $i$ tal que
$k_R(i)=v$ --- não há deduplicação na junção. A condição exige presença nos
dois lados: ``a chave ausente não junta'' (esquerda \cp{sql.c:9279-9285};
direita \cp{sql.c:7358-7363}). A realização distingue os dois modos do dual:
para $=$, corte da árvore (\cp{j_casam}) retorna a fibra inteira; para as
quatro desigualdades, a ordem e intervalos (\cp{sql.c:7713-7717,
9240-9610}) retornam exatamente as fibras relevantes.
\end{definicao}

A \emph{aus\^encia} entra nas jun\c{c}\~oes externas, e a forma como entra \'e
por \emph{fibras vazias} (\cp{sql.c:7679-7689}). Dada a jun\c{c}\~ao interna $P =
R\Join S$, seja $\bot_S$ a linha sem nenhum valor em $J_S$ (a aus\^encia
completa) e $\bot_R$ a sua sim\'etrica:
\[
R\Join^{L,\theta} S = P\ \cup\
\big\{\, r_i\oplus\bot_S : i\in I_R,\ \mathrm{pres}_R(i,a)=1,\
        \nexists\,t : r_i(a)\,\theta\, s_t(b) \,\big\},
\]
\[
R\Join^{R,\theta} S = P\ \cup\
\big\{\, \bot_R\oplus s_t : t\in I_S,\ \mathrm{pres}_S(t,b)=1,\
        \nexists\,i : r_i(a)\,\theta\, s_t(b) \,\big\},
\]
e \cp{FULL OUTER} \'e a uni\~ao das duas extens\~oes
(\cp{j_left = j_right = 1}, \cp{sql.c:7699-7702}). Nota: uma linha com a chave
ausente no seu pr\'oprio lado n\~ao entra em $P$ \emph{e} n\~ao tem fantasma ---
em \cp{LEFT}, a linha de $R$ sem \cp{pres}(i,a) desaparece por completo
(\cp{sql.c:9279-9285}).

\begin{observacao}[O conflito do zero fantasma]
\label{obs:ghost}
\textbf{Interpreta\c{c}\~ao matem\'atica parcial.} Os fantasmas de \cp{LEFT},
\cp{RIGHT} e \cp{FULL} imprimem ``0'' nas c\'elulas do lado ausente e escrevem
no \cp{SqlOut} a c\'elula ``0'' \emph{com o bit de nulo desligado}
(\cp{sql.c:9451-9455, 9510-9512}); o texto diz que o zero \'e ``a
aus\^encia: $b=0$'' (\cp{sql.c:7648, 9451}). Matem\'aticamente queremos uma
\emph{c\'elula ausente}; o canal entrega um \cp{INT4} $0$ \emph{presente}. Do
c\'odigo e do protocolo n\~ao se tira qual das duas leituras \'e a do contrato;
\emph{n\~ao se resolve por suposi\c{c}\~ao}. \stipa{} para as formas externas.
\end{observacao}

\textbf{Contrato implementado.} Os limites s\~ao recusas com o n\'umero, nunca
trunca\c{c}\~ao: o lado esquerdo materializado cabe em \cp{J\_MAXLIN} $=512$
linhas e o total em \cp{J\_MAXCOL} $=\cp{COL\_MAX}=1024$ colunas
(\cp{sql.c:7215-7223, 9254-9276, 9308-9324, 9317-9323}). Uma segunda
jun\c{c}\~ao (tr\^es tabelas) \'e a mesma opera\c{c}\~ao em sequ\^encia: o par de
$R\Join S$ \'e a esquerda de $\Join T$, com o segundo \cp{ON} resolvido contra a
concatena\c{c}\~ao (\cp{sql.c:9532-9604}); n\~ao existe jun\c{c}\~ao tern\'aria.
\stni{} \cp{CROSS}, \cp{NATURAL}, \cp{ON} com mais de uma coluna por lado, ou
constantes no predicado.

\textbf{Status:} jun\c{c}\~ao interna ($\theta$ qualquer, incl. a de tr\^es
tabelas): \sti. \cp{LEFT}/\cp{RIGHT}/\cp{FULL OUTER}: \stipa{} (Obs.~\ref{obs:ghost}).

\subsection{Fibras de chave}
\label{sec:fibras}

Seja $I'=\{i\in I:\operatorname{pres}(i,c)=1\}$ o conjunto de instâncias
vivas com a chave presente na coluna $c$. Defina separadamente
$I_\bot = I\setminus I'$ para as instâncias cuja chave está ausente. O
tratamento de $I_\bot$ é específico de cada operação e não altera a
definição de fibra ordinária.

Para uma função de chave
\[
k:I'\to V,
\]
a \emph{fibra} ordinária de $v$ é
\[
\boxed{
\operatorname{Fib}(k,v)
=
\{i\in I':k(i)=v\}.
}
\]
Trata-se de um conjunto de \emph{índices de instâncias}, logo preserva
duplicadas de conteúdo. As fibras $\{\operatorname{Fib}(k,v):v\in\operatorname{im}(k)\}$
formam uma partição de $I'$.

\subsubsection{Representação por endereçamento recursivo}
\label{sec:fibras-repr}
A implementação indexa as chaves pela \emph{aranha}. Para valores de 32 bits,
\[
B=\mathbb Z/16\mathbb Z,\qquad n=8,\qquad X^{(n)}=B\times X^{(n-1)}.
\]
Tem-se
\[
I'
\xrightarrow{k}
V
\xrightarrow{\operatorname{code}}
X^{(n)}
\xrightarrow{F}
\mathcal P(I'),
\qquad
F(x_v)=\operatorname{Fib}(k,v),\; x_v=\operatorname{code}(v)\in X^{(n)}.
\]
Aqui $X^{(n)}$ é o \emph{carrier de endereçamento} das chaves; $F$ é a
camada relacional de fibras. A \emph{aranha} é a representação/indexação
eficiente de $F$. A coordenada $C^{(n)}$ intervém apenas quando necessário
para descrever intervalos de endereços (desigualdades), não sendo parte
da definição básica de $\operatorname{Fib}(k,v)$.

\subsection{Escritas}
\label{sec:esc}
\begin{itemize}
\item \textbf{\cp{INSERT}} \'e a adi\c{c}\~ao ao multiconjunto:
  \[
  I' = I\,\uplus\,\{\text{linha nova}\},
  \]
  \emph{escrever faz existir}: o valor posto acende $\mathrm{pres}$ e o bit do
  vivo (\cp{insere}, \cp{sql.c:3592}; \cp{insere_muitas}, \cp{sql.c:2967}). Um
  \cp{NULL} expl\'icito grava aus\^encia (\cp{sql.c:3689-3690}); uma coluna n\~ao
  nomeada nasce ausente a menos que um \cp{DEFAULT} a ocupe (\cp{sql.c:3795-3801});
  o envelope recusa o que n\~ao cabe (\cp{sql.c:3665-3690}). O
  \cp{SELECT <constante>} da Sec.~\ref{sec:proj} n\~ao toca o estado.
  \textbf{Status: \sti} (linha \'unica). O multi-\cp{VALUES} corre dentro de uma
  transac\c{c}\~ao de di\'ario (\cp{sql_tx_abre}, \cp{sql.c:1519; 1238-1309});
  o comportamento exacto quando o di\'ario enche a meio de um \cp{VALUES} com
  v\'arias linhas \'e \stni{} neste paper.
\item \textbf{\cp{UPDATE} \cp{SET}} \'e a escrita pontual: $r_i(j)$ troca de
  valor, acende $\mathrm{pres}(i,j)$ quando havia aus\^encia, e
  \cp{SET col = NULL} desliga-a para o dual (\cp{ACAO_SET}, \cp{sql.c:7817-7930,
  6126-6133}). Antes de gravar, revalidam-se \cp{CHECK} e \cp{NOT NULL}/\cp{UNIQUE}
  (\cp{sql.c:9111-9145}); a recusa \"e at\'omica --- o estado n\~ao muda. Um
  racional com denominador $\ne1$ \'e recusado em voz alta
  (\cp{sql.c:7918-7929}).
  \textbf{Status: \sti}
\item \textbf{\cp{DELETE}} \'e a subtrac\c{c}\~ao do multiconjunto: $I'=I\setminus
  \{i\}$ via o bit do vivo (\cp{ACAO_APAGA}, \cp{sql.c:6116, 6303}). Quando a
  linha apontada \'e chave de uma filha, aplica-se a tripla resposta a uma seta
  sem destino --- \cp{RESTRICT} recusa (nada muda), \cp{CASCADE} leva a fibra
  $\pi^{-1}$ atr\'as, \cp{SET NULL} solta-a para o dual (\cp{sql.c:6194-6226}),
  e com v\'arias filhas pergunta-se primeiro a todas se alguma recusa, e s\'o
  depois se age.
  \textbf{Status: \sti}
\item \textbf{Transac\c{c}\~ao e lote.} O ficheiro de sess\~ao tem um di\'ario
  (\cp{sql.c:1238-1309}); o motor aceita v\'arias instru\c{c}\~oes num texto,
  separadas por ``;'', com coment\'arios (batch, \cp{sql.c:16576-16626}).
  \emph{\stco{} que a porta C exponha \cp{COMMIT}/\cp{ROLLBACK} como
  primitivas:} quem as tem \'e o servidor de sess\~ao, n\~ao este caminho.
\end{itemize}

\subsection{\cp{UNION}, vistas, subconsulta}
\label{sec:uni}
\begin{itemize}
\item \textbf{\cp{UNION}} e \cp{UNION ALL} admitem exactamente duas consultas e
  exigem o mesmo n\'umero de colunas. \cp{UNION ALL} \'e a soma de
  multiconjuntos das duas sa\'idas. O \cp{UNION} deduplica \emph{por texto}: o
  crit\'erio \'e a igualdade das c\'elulas formatadas (o \cp{strcmp} sobre a
  linha de texto do \cp{SqlOut}), n\~ao a igualdade dos objectos; duas linhas que
  se l\^em igual s\~ao a mesma (\cp{sql.c:16486-16695}). \textbf{Status: \sti}
  com a particularidade textual declarada.
\item \textbf{Vistas} guardam o texto da consulta; o que uma vista comp\~oe \'e
  exactamente o que o parser do cliente comp\~oe quando a l\^e --- a
  composi\c{c}\~ao \'e textual (\cp{CREATE VIEW}, \cp{sql.c:15539-15596;
  636-641}). \textbf{Status: \sti} para o mecanismo; o fecho alg\'ebrico das
  vistas compos\'iveis \'e \stni{}.
\item \textbf{Subconsulta} na cl\'ausula: s\'o a forma
  ``\cp{(SELECT <col> FROM <t> WHERE …})'' dentro de \cp{IN}/\cp{NOT IN}, e
  reescreve-se para a lista de perten\c{c}a sobre a \'arvore; a fibra vazia da
  subconsulta d\'a uma condi\c{c}\~ao que \emph{nunca casa}
  (\cp{subconsulta_reescreve}, \cp{sql.c:16905; 8848-8919}). Uma consulta n\~ao
  pode trocar a tabela da sess\~ao por baixo de quem a fez: a origem restaura-se
  (\cp{sql.c:8912-8918}). \stni{} subconsultas correlacionadas.
\end{itemize}

\subsection{Delimita\c{c}\~ao do escopo}
\label{sec:escopo}
Fora desta \'algebra --- \stni{} como primitivas do motor --- ficam: \cp{CASE},
janelas/\cp{OVER}, \cp{WITH}/\cp{CTE}, \cp{HAVING} geral, \cp{GROUP BY} com
express\~ao, \cp{ORDER BY} textual, \cp{DISTINCT} multi-coluna, \cp{CROSS}/
\cp{NATURAL} e \cp{ON} com mais de uma coluna por lado, subconsulta
correlacionada, \cp{INSERT..SELECT}, \cp{UPDATE..FROM}/\cp{JOIN},
\cp{DELETE..USING}, \cp{SELECT count(*)} sem \cp{FROM}, e a l\'ogica tern\'aria.
O motor aceita ainda um punhado de opera\c{c}\~oes n\~ao-padr\~ao ---
\cp{ACUMULA}/\cp{DELTA}, o bloco matricial (\cp{triade}, \cp{edo},
\cp{global}, \cp{completa}, \dots), \cp{BUSCA}/\cp{ACHA}/\cp{DISTANCIA},
\cp{GET}/\cp{PARSE}/\cp{MINERA} e os escapes de shell \cp{BASH}/\cp{POWERSHELL}/
\cp{NODE} (\cp{sql.c:9794-11600, 15119, 15800-15802}). Elas s\~ao
gramaticalmente aceites, mas n\~ao s\~ao modeladas aqui: \stco{} nesta sec\c{c}\~ao
(o papel que teriam numa \'algebra do Tiffany est\'a por escrever).

%=======================================================================
\section{N\~ao estabelecido pela implementa\c{c}\~ao atual}
\label{sec:naoest}

Cada item a seguir \'e \textbf{N\~ao estabelecido pela implementa\c{c}\~ao
atual:} — o paper n\~ao o afirma, e o leitor n\~ao deve l\^e-lo no c\'odigo.

\begin{enumerate}
\item \textbf{Probabilidade $p_q$.} N\~ao existe divis\~ao pelo cardinal:
  n\~ao se calcula $p_q(j)=G_q(j)/\lvert I\rvert$ nem se emite. A sa\'ida \'e
  a contagem inteira $G_q$ e o denominador $\lvert I\rvert$ na coluna
  \cp{i_total}.
\item \textbf{Funcional $L(p_q)$.} N\~ao \'e calculado no caminho SQL;
  \cp{jev_escore} (que devolve o par $(N,\lvert I\rvert)$ com
  $N=\sum_i s_i G_q(s_i)$) existe na porta C mas \cp{sql_hist_jev} n\~ao o
  chama. A leitura de $N/\lvert I\rvert$ \'e, no contrato do \cp{jev_api.h},
  \emph{externa} ao motor.
\item \textbf{Realiza\c{c}\~ao armazenada.} $\pi$ n\~ao \'e guardada: o
  c\'odigo percorre as linhas vivas, conta por valor no tamp\~ao e deita fora
  as linhas. S\'o o campo $G$ entra na arena. Distribui\c{c}\~oes conjuntas
  entre duas colunas n\~ao s\~ao constru\'idas aqui.
\item \textbf{Colunas fantasma sem nomes.} Numa tabela sem nomes guardados,
  \cp{col_indice} cai na letra (de $a..z$ a \'indices $0..25$) sem validar
  $\cp{oc}<\cp{nc}$ (\cp{sql.c:2075}); com $c$ colunas e $\cp{nc}\le25$, G2/G3
  podem ent\~ao ler presen\c{c}a e valor de uma c\'elula fora de $J$. O
  contracto n\~ao exclui este caso.
\item \textbf{Corpo largo e o limite de colunas.} \cp{col_indice} s\'o procura
  em $\cp{S\_COLNOME\_N}=1024$ colunas; para $j\ge\cp{COL\_MAX}$,
  \cp{corpo\_de} devolve $\{0,0\}$ cujo \cp{total}=0 \'e \cp{CORPO\_INTEIRO}
  (\cp{sql.c:1715-1719}) — G1 passaria a ver uma coluna fantasma vazia. A
  recusa directa vem de E2/E3, n\~ao de G1.
\item \textbf{Marginal $d$-dimensional.} A ponte \cp{jev_marginal_d} (o
  \cp{sec:byte-carry} do \cp{jev.tex}) n\~ao \'e usada neste caminho: as
  projec\c{c}\~oes aqui s\~ao \'indices de intervalo, n\~ao coordenadas de
  $\mathcal{X}_d$.
\item \textbf{Estado global.} \cp{jev_hist_erros} e \cp{cels\_fora} s\~ao
  contadores \cp{static}, n\~ao observ\'aveis pela API; o papel que teriam
  numa garantia externa n\~ao est\'a escrito em lado nenhum.
\item \textbf{Conserva\c{c}\~ao como prova independente.} A guarda final soma
  a sa\'ida da pr\'opria primitiva; n\~ao \'e uma segunda implementa\c{c}\~ao.
  Uma margem incorrecta e \emph{consistentemente} auto-somada n\~ao seria
  apanhada aqui (ver Obs. da Sec.~\ref{sec:nivel1}).
\end{enumerate}

%=======================================================================
\section{Tabela de correspond\^encias}
\label{sec:corr}

\begin{center}
\begin{tabular}{@{}p{4.2cm} p{4.0cm} p{5.4cm} p{2.8cm}@{}}
\toprule
\textbf{opera\c{c}\~ao SQL} &
\textbf{defini\c{c}\~ao matem\'atica} &
\textbf{c\'odigo C} &
\textbf{status} \\
\midrule
\cp{SELECT} (projec\c{c}\~ao, express\~oes) &
$\pi_p$, tensor $a+b$ e s\'eries exp/sin/cos/log &
\cp{varre} 7751--7810, emiss\~ao 12902--13164 &
\sti \\
\cp{WHERE} (AND/OR/NOT, =,<,<=,>,>=, BETWEEN/IN, IS NULL) &
$\sigma_\varphi$ com $C(\varphi)$ e aus\^encia $=$ fila ausente &
\cp{le_where}/\cp{le_expr} $\sim$4922--4810, 8933--8995 &
\sti \\
\cp{DISTINCT} (1 coluna) &
imagem por fibra (classe + aus\^encia) &
\cp{sql.c:12845--12900} &
\sti \\
\cp{JOIN INNER} (\cp{a.x $\theta$ b.y}, 1:1, $\theta\in\{=,<,\le,>,\ge\}$) &
$R\Join^\theta S$ por pares, presen\c{c}a nos dois lados, corte/ordem &
\cp{le_join} 7675--7739, \cp{j_carrega_direita} 7338, materializa\c{c}\~ao 9240--9610 &
\sti \\
\cp{JOIN LEFT} &
$R\Join^L S$; coluna fantasma ``0'' presente no \cp{SqlOut} &
\cp{sql.c:7680--7690, 9440--9461} &
\stipa \\
\cp{JOIN RIGHT/FULL OUTER} &
formas externas; mesmo fantasma para o lado ausente &
\cp{sql.c:7695--7702, 9502--9528} &
\stipa \\
\cp{GROUP BY} (1--2 colunas) &
$\bigsqcup_k R_k$, $k=g(i)$; aus\^encia = grupo pr\'oprio; \cp{count} sempre emitido &
\cp{sql.c:8362--8390, 12612--12836} &
\sti \\
\cp{SUM/MAX/MIN/AVG} (num\'ericas) &
sobre fibras (soma/\'extremo); \cp{AVG} racional reduzido; ausentes n\~ao entram &
\cp{sql.c:12527--12610} &
\sti \\
\cp{COUNT(*)} &
$|I(\varphi)|$ (popcount do campo) &
\cp{sql.c:15929--15935, 8664--8679} &
\sti \\
\cp{COUNT(DISTINCT)} (1 coluna) &
n\'umero de classes (inclui aus\^encia) &
\cp{sql.c:9000--9044} &
\sti \\
\cp{HAVING} (s\'o \cp{count(*)} $\theta n$) &
filtragem de $|R_k|$ &
\cp{sql.c:8393--8429, 12749--12753} &
\sti \\
\cp{ORDER BY} (1--2 colunas, ASC/DESC) &
permuta\c{c}\~ao por r\'egua do corpo (declara norma); ausentes ao fim; texto recusado &
\cp{sql.c:8431--8500, 9617--9701} &
\sti \\
\cp{LIMIT/OFFSET} &
fatia $[o, o+\ell)$ (OFFSET antes de LIMIT) &
\cp{sql.c:8504--8531, 12906--12910} &
\sti \\
\cp{INSERT} (uma linha; multi-\cp{VALUES} em TX) &
$I'=I\uplus\{r\}$; \cp{NULL}$=$aus\^encia; escrever acende \cp{S_PRES}; envelope recusa &
\cp{insere_muitas} 2967, \cp{insere} 3592, 3681--3692 &
\sti \\
\cp{UPDATE SET} &
escrita pontual; recusa racional den!=1; revalida \cp{CHECK}/\cp{NOT NULL}/\cp{UNIQUE} &
\cp{sql.c:7817--7930, 9111--9145, 6126--6133} &
\sti \\
\cp{DELETE} (FK \cp{RESTRICT}/\cp{CASCADE}/\cp{SET NULL}) &
$I'=I\setminus\{i\}$ (bit vivo); cascata/solta com regra de duas passagens &
\cp{sql.c:6116, 6303, 6194--6226} &
\sti \\
\cp{UNION}/\cp{UNION ALL} (2 consultas, ncols iguais) &
soma de multiconjuntos; \cp{UNION} dedup por \cp{strcmp} textual das linhas &
\cp{sql.c:16486--16695} &
\sti \\
\cp{CREATE VIEW} &
composi\c{c}\~ao textual (guardado como texto) &
\cp{sql.c:15539--15596} &
\sti \\
\cp{CREATE INDEX [UNIQUE]} ON t(c); \cp{CREATE TYPE ENUM}; \cp{ALTER/DROP} &
metadados/\'indice na \'arvore; limita\c{c}\~oes documentadas &
\cp{sql.c:15452--15727} &
\stco \\
\cp{EXISTS} (forma restrita: \cp{SELECT * FROM t} ou \cp{col FROM t}) &
quantificador constante para a consulta; \cp{NOT EXISTS} complemento &
\cp{sql.c:8280--8331, 8982--8986} &
\stco \\
\cp{subconsulta \cp{(SELECT)} em \cp{IN}/\cp{NOT IN}} (s\'o lista) &
reescrita para perten\c{c}a sobre \'arvore; fibra vazia $=$ nunca casa; restaura tabela &
\cp{subconsulta_reescreve} 16905, 8848--8919 &
\stco \\
\cp{SELECT <constante>} (sem \cp{FROM}) &
express\~ao sem vari\'aveis (teste de liga\c{c}\~ao) &
\cp{sql.c:15899--15927} &
\stco \\
Opera\c{c}\~oes n\~ao-padr\~ao (mat\_op, ACUMULA/DELTA, BUSCA/ACHA/MINERA, escapes shell) &
n\~ao modeladas nesta \'algebra do Tiffany & citado &
\stco \\
\cp{CROSS}, \cp{NATURAL}, \cp{ON} multi-coluna, \cp{CASE}, \cp{OVER}, \cp{WITH}, \cp{HAVING} geral, \cp{INSERT..SELECT}, \cp{UPDATE..FROM}, \cp{DELETE..USING}, subconsulta correlacionada, \cp{DISTINCT} multi-coluna, \cp{GROUP BY} com express\~ao &
n\~ao implementadas / n\~ao estabelecidas & citado &
\stni \\
\bottomrule
\end{tabular}
\end{center}

\medskip
\paragraph{Notas:} ``implementado mas interpreta\c{c}\~ao matem\'atica parcial''
aplica-se \`as jun\c{c}\~oes externas, porque o fantasma ``0'' emitido no
\cp{SqlOut} (\cp{sql.c:9451, 9510}) n\~ao tem o bit de nulo aceso, e o
protocolo n\~ao permite decidir sozinho entre ``valor 0 presente'' e
``c\'elula ausente''. Para \cp{WHERE}, a sem\^antica de aus\^encia/NULL \'e
determinada (bivalente): as linhas que citam coluna ausente s\~ao exclu\'idas
(\cp{sql.c:8933--8941}), \cp{IS [NOT] NULL} decide pelo \cp{S_PRES}, e
\cp{EXISTS} \'e um quantificador. \textbf{COUNT} \'e sempre por campo/popcount
conforme estabelecido; \cp{DISTINCT} \'e de uma coluna s\'o; \cp{JOIN} limitado
a um predicado de uma coluna por lado e a at\'e 3 tabelas em sequ\^encia; o
tamanho da \'arvore e $\cp{SQL\_OUT\_MAX\_*}$ recusa (nunca trunca); a deduplica\c{c}\~ao
do \cp{UNION} \'e por texto (\cp{strcmp}); \cp{OFFSET} antes de \cp{LIMIT};
\cp{ORDER BY} em corpos quadr\'aticos usa a norma declarada; \cp{GROUP BY}
emite sempre \cp{count}; \cp{COUNT(DISTINCT)} conta a classe da aus\^encia;
\cp{SELECT count(*)} sem \cp{FROM} n\~ao existe; \cp{EXISTS} s\'o na forma
restrita. O di\'ario do multi-\cp{VALUES} em transac\c{c}\~ao permanece com o
ponto ``indetermina\c{c}\~ao de estado'' citado no invent\'ario (\stni).

\medskip
E o mapa de recusas (N\'ivel 1 $\leftrightarrow$ N\'ivel 2):
\begin{center}
\begin{tabular}{@{}c l@{}}
\toprule
G7& tabela/coluna inexistentes, tabela vazia \\
G1& corpo != \cp{CORPO\_INTEIRO} (E1) \\
G2& \emph{presen\c{c}a incompleta} (E2) \\
G3& valor fora do tamp\~ao $[0,16384)$ (E3) \\
\cp{spec\_decode}& \emph{spec inv\'alida} (E6) \\
G4& $k\notin[1,16]$ (E4; inalcan\c{c}avel via \cp{spec}) \\
G5& $X<1$ ou $X+2k>16230$ (E5) \\
F17& $v\ge X$ (E3) \\
G6& fronteiras fora de $(0,X)$ ou n\~ao crescentes (E6) \\
\bottomrule
\end{tabular}
\end{center}

%=======================================================================
\section{Compila\c{c}\~ao e proveni\^encia}

\paragraph{Para este paper (auto-contido).}
\[
\cp{cd papers \&\& pdflatex sql.tex}
\]
Usa \cp{\\input\{gkcapa\}} e o estilo de \cp{naturais.tex}; n\~ao precisa de
nada mais (o \cp{pdflatex} basta; duas passagens para o \cp{hyperref}).

\paragraph{Para o c\'odigo (grafo de depend\^encias do motor SQL).}
\[
\cp{gcc -MM -Ibanco -Ilib banco/sql.c banco/pgwire.c}
\]
(\cp{debian/rules} \'e a receita oficial; o \cp{tiffany-db} n\~ao tem
Makefile/CMake.)

\paragraph{Proveni\^encia dos n\'umeros citados.}
\cp{tiffany-db/JEV.md}; \cp{PROPOSTA\_TIFFANY\_DB.md};
\cp{FASE\_4\_9\_R3\_STATUS.md}; \cp{FASE\_4\_9\_R3\_RESULTADO.md}.
O alinhamento de nota\c{c}\~ao ($\pi$, $G$, $q$, $G_q$) est\'a em
\cp{tiffany/redes/jev.tex} (\cp{def:realizacao}, \cp{prop:gq-conserva}) —
usado aqui como vocabul\'ario, n\~ao como justifica\c{c}\~ao.

\end{document}